\documentclass[12pt,a4paper]{article}

% --- Настройка языка и шрифтов для LuaLaTeX ---
\usepackage{fontspec}
\defaultfontfeatures{Ligatures=TeX,Scale=MatchLowercase}

\setmainfont{Alice-Regular}[
    Path=font/,
    Extension=.ttf,
    BoldFont = Alice-Regular,
    ItalicFont = Alice-Regular,
    BoldItalicFont = Alice-Regular,
    BoldFeatures = {FakeBold=1.6},
    ItalicFeatures = {FakeSlant=0.2},
    BoldItalicFeatures = {FakeBold=1.6,FakeSlant=0.2}
]

\newfontfamily\cyrillicfont{Alice-Regular}[
    Path=font/,
    Extension=.ttf,
    BoldFont = Alice-Regular,
    ItalicFont = Alice-Regular,
    BoldItalicFont = Alice-Regular,
    BoldFeatures = {FakeBold=1.6},
    ItalicFeatures = {FakeSlant=0.2},
    BoldItalicFeatures = {FakeBold=1.6,FakeSlant=0.2}
]

\usepackage{polyglossia}
\setmainlanguage{russian}
\setsansfont{TeX Gyre Heros}
\setmonofont{TeX Gyre Cursor}

\usepackage{graphicx}        					          % Графика
\usepackage[protrusion=false]{microtype}       	% Микротипографика
\usepackage{amsmath, amsfonts, amssymb, amsthm}	% Математика
\usepackage{braket} 							              % Braket нотация для квантовой механики
\usepackage{unicode-math} 						          % Современный пакет для математики в LuaLaTeX
\usepackage{longtable}       					          % Таблицы с разрывом страниц
\usepackage{tikz}           					          % Рисование
\usepackage[europeanresistors]{circuitikz}		  % Рисование электрических схем в европейском стиле
\usepackage{listings}							              % Для отображения кода
\usepackage{tabularx}       					          % Для таблиц во всю ширину
\usepackage{array}		  						            % Для расширенных возможностей таблиц
\usepackage{float}          					          % Для точного позиционирования фигур
\usepackage{pgfplots}                           % Для построения графиков
\usepackage{caption}                            % Для настройки подписей к рисункам и таблицам
\usepackage{xfp}                                % Для вычислений в LaTeX
\pgfplotsset{compat=1.18}
\usepackage[a4paper,left=20mm,right=20mm,top=20mm,bottom=20mm]{geometry}

\usetikzlibrary{arrows.meta}

% --- Колонтитулы ---
\usepackage{fancyhdr}
\pagestyle{fancy}
\fancyhf{}

\renewcommand{\headrulewidth}{0.4pt}
\renewcommand{\footrulewidth}{0.4pt}
\renewcommand\lstlistingname{Исходный код}

\lstdefinestyle{main}{
	backgroundcolor=\color{gray!8},   
	commentstyle=\color{green!50!black},
	keywordstyle=\color{blue!80!black},
	numberstyle=\small\color{darkgray},
	stringstyle=\rmfamily\color{orange!90!black},
	basicstyle=\ttfamily\small,
	breakatwhitespace=false,         
	breaklines=true,                   
	keepspaces=true,                 
	numbers=left,                    
	numbersep=8pt,                  
	showspaces=false,                
	showstringspaces=false,
	showtabs=false,                  
	tabsize=4
}	

\lstset{style=main}

\newcolumntype{C}{>{\centering\arraybackslash}X}
\renewcommand{\arraystretch}{1.25}

% =========================================================
%         СЛУЖЕБНЫЕ КОМАНДЫ ДЛЯ ВЫЧИСЛЕНИЙ
% =========================================================
\newcommand{\CalcTwo}  [1]{\fpeval{round((#1),2)}}
\newcommand{\CalcThree}[1]{\fpeval{round((#1),3)}}
\newcommand{\CalcFour} [1]{\fpeval{round((#1),4)}}
\newcommand{\CalcFive} [1]{\fpeval{round((#1),5)}}

% ---------------------------------------------------------
% ОБЩИЕ ДАННЫЕ
% ---------------------------------------------------------
\def\VariantN{1}          					         % Номер варианта
\def\StudentName{Андреев Иван Алексеевич}  	 % ФИО студента
\def\TeacherName{NOTHING} % ФИО преподавателя
\def\LabNumber{11}							               % Номер лабораторной работы
\def\Discipline{Атомная физика}				           % Название дисциплины
\def\LabTitle{Определение удельного заряда электрона}	 % Название лабораторной работы

\def\dOne{\fpeval{10 * 10^(-2)}} % см -> м
\def\dTwo{\fpeval{8 * 10^(-2)}} % см -> м
\def\dThree{\fpeval{6 * 10^(-2)}} % см -> м
\def\dFour{\fpeval{4 * 10^(-2)}} % см -> м

\def\Ione{\fpeval{1.47}} % А
\def\Itwo{\fpeval{1.87}} % А
\def\Ithree{\fpeval{2.44}} % А
\def\Ifour{\fpeval{3.86}} % А

\def\U{\fpeval{240}} % В

\title{
	МИНОБРНАУКИ РОССИИ

	федеральное государственное автономное образовательное учреждение

	высшего образования

	«Национальный исследовательский университет
	«Московский институт электронной техники»
	\vspace{2em}
}
\author{
}
\date{}

\begin{document}

\maketitle
\thispagestyle{fancy}

\begin{center}
\Large Отчёт по Лабораторной работе $\LabNumber$\\
по дисциплине «\Discipline»\\
«\LabTitle»
\end{center}

\begin{center}
\Large Вариант \VariantN
\end{center}

\vspace{12em}

\begin{table}[h]
\noindent
\begin{tabular*}{\textwidth}{@{\extracolsep{\fill}} p{0.48\textwidth} p{0.52\textwidth}}
Студент, группа & \StudentName, ЭН-$25$ \\
[0.5em] Руководитель & \TeacherName
\end{tabular*}
\end{table}

\cfoot{Москва 2026}

\pagebreak

\fancyhead[L]{\LabTitle} 
\fancyhead[R]{\Discipline}
\fancyfoot[C]{\thepage}

Способ 1

\def\BxOne{\fpeval{6.9184 * 10^(-4) * \Ione}} % Тл
\def\BxTwo{\fpeval{6.9184 * 10^(-4) * \Itwo}} % Тл
\def\BxThree{\fpeval{6.9184 * 10^(-4) * \Ithree}} % Тл
\def\BxFour{\fpeval{6.9184 * 10^(-4) * \Ifour}} % Тл

\begin{equation*}
	\begin{aligned}
		B_\text{x1} &= 6.9184 \cdot 10^{-4} \cdot I_1 = \CalcFive{6.9184 * 10^(-4) * \Ione * 10^{3}} \cdot 10^{-3}  \text{ Тл} \\
		B_\text{x2} &= 6.9184 \cdot 10^{-4} \cdot I_2 = \CalcFive{6.9184 * 10^(-4) * \Itwo * 10^{3}} \cdot 10^{-3} \text{ Тл} \\
		B_\text{x3} &= 6.9184 \cdot 10^{-4} \cdot I_3 = \CalcFive{6.9184 * 10^(-4) * \Ithree * 10^{3}} \cdot 10^{-3} \text{ Тл} \\
		B_\text{x4} &= 6.9184 \cdot 10^{-4} \cdot I_4 = \CalcFive{6.9184 * 10^(-4) * \Ifour * 10^{3}} \cdot 10^{-3} \text{ Тл}
	\end{aligned}
\end{equation*}

\def\omegaOne{\fpeval{8 * \U / (\dOne^2 * \BxOne^2)}} % Кл/кг
\def\omegaTwo{\fpeval{8 * \U / (\dTwo^2 * \BxTwo^2)}} % Кл/кг
\def\omegaThree{\fpeval{8 * \U / (\dThree^2 * \BxThree^2)}} % Кл/кг
\def\omegaFour{\fpeval{8 * \U / (\dFour^2 * \BxFour^2)}} % Кл/кг

\begin{equation*}
	\begin{aligned}
		\omega_\text{1} &= \frac{8U_1}{d^2B_1^2} = \CalcFive{8 * \U / (\dOne^2 * \BxOne^2) * 10^{-11}} \cdot 10^{11} \text{ Кл/кг} \\
		\omega_\text{2} &= \frac{8U_2}{d^2B_2^2} = \CalcFive{8 * \U / (\dTwo^2 * \BxTwo^2) * 10^{-11}} \cdot 10^{11} \text{ Кл/кг} \\
		\omega_\text{3} &= \frac{8U_3}{d^2B_3^2} = \CalcFive{8 * \U / (\dThree^2 * \BxThree^2) * 10^{-11}} \cdot 10^{11} \text{ Кл/кг} \\
		\omega_\text{4} &= \frac{8U_4}{d^2B_4^2} = \CalcFive{8 * \U / (\dFour^2 * \BxFour^2) * 10^{-11}} \cdot 10^{11} \text{ Кл/кг}
	\end{aligned}
\end{equation*}

\begin{equation*}
	\omega_\text{т} = \frac{e}{m_e} = \CalcFive{1.6 * 10^{-19} / (9.1 * 10^(-31)) * 10^{-11}} \cdot 10^{11} \text{ Кл/кг}
\end{equation*}

\def\deltaU{\fpeval{0.01 + 5 / \U}}
\def\deltaIone{\fpeval{0.02 + 5 * 0.01 / \Ione}}
\def\deltaItwo{\fpeval{0.02 + 5 * 0.01 / \Itwo}}
\def\deltaIthree{\fpeval{0.02 + 5 * 0.01 / \Ithree}}
\def\deltaIfour{\fpeval{0.02 + 5 * 0.01 / \Ifour}}

\begin{equation*}
	\delta U = 0.01 + \frac{5}{U} = \CalcTwo{(0.01 + 5/\U) * 100} \%
\end{equation*}

\begin{equation*}
	\begin{aligned}
		\delta I_1 &= 0.02 + \frac{5 \cdot 0.01}{I_1} = \CalcTwo{(0.02 + 5 * 0.01/\Ione) * 100} \% \\
		\delta I_2 &= 0.02 + \frac{5 \cdot 0.01}{I_2} = \CalcTwo{(0.02 + 5 * 0.01/\Itwo) * 100} \% \\
		\delta I_3 &= 0.02 + \frac{5 \cdot 0.01}{I_3} = \CalcTwo{(0.02 + 5 * 0.01/\Ithree) * 100} \% \\
		\delta I_4 &= 0.02 + \frac{5 \cdot 0.01}{I_4} = \CalcTwo{(0.02 + 5 * 0.01/\Ifour) * 100} \%
	\end{aligned}
\end{equation*}

\begin{equation*}
	\begin{aligned}
		\Delta \omega_1 &= \omega_1 \cdot \left(\delta U + 2 \delta I_1\right) = \CalcFive{\omegaOne * (\deltaU + 2 * \deltaIone) * 10^{-11}} \cdot 10^{11} \text{ Кл/кг} \\
		\Delta \omega_2 &= \omega_2 \cdot \left(\delta U + 2 \delta I_2\right) = \CalcFive{\omegaTwo * (\deltaU + 2 * \deltaItwo) * 10^{-11}} \cdot 10^{11} \text{ Кл/кг} \\
		\Delta \omega_3 &= \omega_3 \cdot \left(\delta U + 2 \delta I_3\right) = \CalcFive{\omegaThree * (\deltaU + 2 * \deltaIthree) * 10^{-11}} \cdot 10^{11} \text{ Кл/кг} \\
		\Delta \omega_4 &= \omega_4 \cdot \left(\delta U + 2 \delta I_4\right) = \CalcFive{\omegaFour * (\deltaU + 2 * \deltaIfour) * 10^{-11}} \cdot 10^{11} \text{ Кл/кг}
	\end{aligned}
\end{equation*}

Способ 2

\def\OneDividedOneSquared{\fpeval{1/\dOne^2}} % м^{-2}
\def\OneDividedTwoSquared{\fpeval{1/\dTwo^2}} % м^{-2}
\def\OneDividedThreeSquared{\fpeval{1/\dThree^2}} % м^{-2}
\def\OneDividedFourSquared{\fpeval{1/\dFour^2}} % м^{-2}

\def\BxOneSquared{\fpeval{\BxOne^2}} % Тл^2
\def\BxTwoSquared{\fpeval{\BxTwo^2}}
\def\BxThreeSquared{\fpeval{\BxThree^2}}
\def\BxFourSquared{\fpeval{\BxFour^2}}

\begin{tabularx}{\textwidth}{|C|C|C|C|C|C|}
	\hline
	\textbf{п/п} & $d$\textbf{, см} & $\OneDividedOneSquared$, м$^{-2}$ & $I$\textbf{, А} & $B_x$\textbf{, Тл} & $B_x^2$\textbf{, Тл$^2$} \\
	\hline
	$1$ & $\fpeval{\dOne * 10^{2}}$ & $\CalcTwo{\OneDividedOneSquared}$ & $\Ione$ & $\CalcTwo{\BxOne * 10^{3}} \cdot 10^{-3}$ & $\CalcTwo{\BxOneSquared * 10^{6}} \cdot 10^{-6}$  \\
	\hline
	$2$ & $\fpeval{\dTwo * 10^{2}}$ & $\CalcTwo{\OneDividedTwoSquared}$ & $\Itwo$ & $\CalcTwo{\BxTwo * 10^{3}} \cdot 10^{-3}$ & $\CalcTwo{\BxTwoSquared * 10^{6}} \cdot 10^{-6}$ \\
	\hline
	$3$ & $\fpeval{\dThree * 10^{2}}$ & $\CalcTwo{\OneDividedThreeSquared}$ & $\Ithree$ & $\CalcTwo{\BxThree * 10^{3}} \cdot 10^{-3}$ & $\CalcTwo{\BxThreeSquared * 10^{6}} \cdot 10^{-6}$ \\
	\hline
	$4$ & $\fpeval{\dFour * 10^{2}}$ & $\CalcTwo{\OneDividedFourSquared}$ & $\Ifour$ & $\CalcTwo{\BxFour * 10^{3}} \cdot 10^{-3}$ & $\CalcTwo{\BxFourSquared * 10^{6}} \cdot 10^{-6}$ \\
	\hline
\end{tabularx}

\def\DeltaBxOneSquared{\fpeval{\BxOne^2 * 2 * \deltaIone}} % Тл^2
\def\DeltaBxTwoSquared{\fpeval{\BxTwo^2 * 2 * \deltaItwo}} % Тл^2
\def\DeltaBxThreeSquared{\fpeval{\BxThree^2 * 2 * \deltaIthree}} % Тл^2
\def\DeltaBxFourSquared{\fpeval{\BxFour^2 * 2 * \deltaIfour}} % Тл^2


\begin{equation*}
	\begin{aligned}
		\Delta B_{x1}^2 &= B_{x1}^2 \cdot 2 \delta I_1 = \CalcTwo{\DeltaBxOneSquared * 10^{6}} \cdot 10^{-6} \text{ Тл} \\
		\Delta B_{x2}^2 &= B_{x2}^2 \cdot 2 \delta I_2 = \CalcTwo{\DeltaBxTwoSquared * 10^{6}} \cdot 10^{-6} \text{ Тл} \\
		\Delta B_{x3}^2 &= B_{x3}^2 \cdot 2 \delta I_3 = \CalcTwo{\DeltaBxThreeSquared * 10^{6}} \cdot 10^{-6} \text{ Тл} \\
		\Delta B_{x4}^2 &= B_{x4}^2 \cdot 2 \delta I_4 = \CalcTwo{\DeltaBxFourSquared * 10^{6}} \cdot 10^{-6} \text{ Тл}
	\end{aligned}
\end{equation*}

\begin{equation*}
	\gamma_\text{max} = \frac{\fpeval{1/\dOne^2} - \fpeval{1/\dFour^2}}
							 {(\CalcTwo{\BxOne^2 * 10^{6}} \cdot 10^{-6} + 
							 \CalcTwo{\BxOne^2 * 2 * \deltaIone * 10^{6}} \cdot 10^{-6}) 
							 - (\CalcTwo{\BxFour^2 * 10^{6}} \cdot 10^{-6} - 
							 \CalcTwo{\BxFour^2 * 2 * \deltaIfour * 10^{6}} \cdot 10^{-6})} = 
							 \CalcTwo{((1/\dOne^2) - (1/\dFour^2)) / ((\BxOne^2 + \BxOne^2 * 2 * \deltaIone) - (\BxFour^2 - \BxFour^2 * 2 * \deltaIfour))} \text{ м}^{-2}\text{Тл}^{-2}
\end{equation*}

\begin{equation*}
	\gamma_\text{min} = \frac{\fpeval{1/\dOne^2} - \fpeval{1/\dFour^2}}
							 {(\CalcTwo{\BxOne^2 * 10^{6}} \cdot 10^{-6} - 
							 \CalcTwo{\BxOne^2 * 2 * \deltaIone * 10^{6}} \cdot 10^{-6}) 
							 - (\CalcTwo{\BxFour^2 * 10^{6}} \cdot 10^{-6} + 
							 \CalcTwo{\BxFour^2 * 2 * \deltaIfour * 10^{6}} \cdot 10^{-6})} = 
							 \CalcTwo{((1/\dOne^2) - (1/\dFour^2)) / ((\BxOne^2 - \BxOne^2 * 2 * \deltaIone) - (\BxFour^2 + \BxFour^2 * 2 * \deltaIfour))} \text{ м}^{-2}\text{Тл}^{-2}
\end{equation*}

\def\gammaMax{\fpeval{((1/\dOne^2) - (1/\dFour^2)) / ((\BxOne^2 + \BxOne^2 * 2 * \deltaIone) - (\BxFour^2 - \BxFour^2 * 2 * \deltaIfour))}} % м^{-2}Тл^{-2}
\def\gammaMin{\fpeval{((1/\dOne^2) - (1/\dFour^2)) / ((\BxOne^2 - \BxOne^2 * 2 * \deltaIone) - (\BxFour^2 + \BxFour^2 * 2 * \deltaIfour))}} % м^{-2}Тл^{-2}

\begin{equation*}
	\braket{\gamma} = \frac{\gamma_\text{max} + \gamma_\text{min}}{2} = \CalcTwo{(\gammaMax + \gammaMin) / 2} \text{ м}^{-2}\text{Тл}^{-2}
\end{equation*}

\begin{equation*}
	\braket{\Delta\gamma} = \frac{\gamma_\text{max} - \gamma_\text{min}}{2} = \CalcTwo{(\gammaMax - \gammaMin) / 2} \text{ м}^{-2}\text{Тл}^{-2}
\end{equation*}

\begin{equation*}
	\omega = 8U \cdot \braket{\gamma} = \CalcFive{8 * \U * (\gammaMax + \gammaMin) / 2 * 10^{-11}} \cdot 10^{11} \text{ Кл/кг}
\end{equation*}

\begin{equation*}
	\Delta \omega = 8U \cdot \braket{\Delta\gamma} = \CalcFive{8 * \U * (\gammaMax - \gammaMin) / 2 * 10^{-11}} \cdot 10^{11} \text{ Кл/кг}
\end{equation*}

\begin{center}
\begin{tikzpicture}
\begin{axis}[
	width=15cm,
	height=10cm,
	xlabel={$B_x^2,\ \mathrm{T}^2$},
	ylabel={$\dfrac{1}{d^2},\ \mathrm{m}^{-2}$},
	xmin=-3.565e-7,
	xmax=8e-6,
	ymin=-31.25,
	ymax=656.25,
	axis lines=middle,
	xtick={0,1e-6,...,8e-6},
	ytick={0,100,...,700},
	minor x tick num=8,
	minor y tick num=8,
	grid=both,
	major grid style={line width=1pt, draw=black!100},
	minor grid style={line width=0.5pt, draw=black!50},
	xticklabel style={/pgf/number format/fixed,/pgf/number format/precision=1},
	yticklabel style={/pgf/number format/fixed,/pgf/number format/precision=0},
	tick align=outside,
	tick style={black},
]

% Экспериментальные точки
\addplot[
    only marks,
    mark=*,
    mark size=2.5pt
]
coordinates {
    (\BxOneSquared,\OneDividedOneSquared)
    (\BxTwoSquared,\OneDividedTwoSquared)
    (\BxThreeSquared,\OneDividedThreeSquared)
    (\BxFourSquared,\OneDividedFourSquared)
};

\addplot[
    only marks,
    mark=|,
    mark size=5pt
]
coordinates {
    (\BxOneSquared + \DeltaBxOneSquared,\OneDividedOneSquared)
    (\BxTwoSquared + \DeltaBxTwoSquared,\OneDividedTwoSquared)
    (\BxThreeSquared + \DeltaBxThreeSquared,\OneDividedThreeSquared)
    (\BxFourSquared + \DeltaBxFourSquared,\OneDividedFourSquared)
	(\BxOneSquared - \DeltaBxOneSquared,\OneDividedOneSquared)
	(\BxTwoSquared - \DeltaBxTwoSquared,\OneDividedTwoSquared)
	(\BxThreeSquared - \DeltaBxThreeSquared,\OneDividedThreeSquared)
	(\BxFourSquared - \DeltaBxFourSquared,\OneDividedFourSquared)
};

% Линейная аппроксимация по этим четырём точкам:
% y = 8.5689e7 * x + 18.124
\addplot[
    domain=-3.565e-7:8e-6,
    samples=2,
    thick,
    dashed
]
{8.5689e7*x + 18.124};

\end{axis}
\end{tikzpicture}
\end{center}
\end{document}